\subsection{Path way: comparación de ida y vuelta}\label{sec:vuelta}


Cada tabla conserva la numeración TTL desde su propio origen, así que las filas de ida y vuelta no identifican automáticamente el mismo router. Para comparar el orden invierto por separado la secuencia de vuelta. Hay que tener en cuenta que un traceroute informa la interfaz por la que entra el paquete, de modo que, aunque el camino físico fuera idéntico en ambos sentidos, cada sentido mostraría direcciones distintas del mismo equipo (las dos puntas de un enlace). Por eso un índice de Jaccard bajo no equivale por sí solo a una ruta asimétrica, y direcciones distintas no demuestran equipos diferentes. Como no hice una prueba independiente de resolución de alias, me limito a comparar interfaces, sus ASN y su orden.

%se invierte por separado la secuencia de vuelta. Direcciones diferentes no demuestran equipos diferentes y no se identificaron alias de interfaces mediante una prueba independiente.

\subsubsection{Valparaíso}
\begingroup
\fontsize{9}{11}\selectfont
\addtocounter{table}{-1}
\begin{longtable}{p{2.1cm}p{1.7cm}p{2cm}p{2cm}p{2.2cm}p{2.7cm}}
\toprule
\textbf{Protocolo} & \textbf{IP ida r2} & \textbf{IP vuelta r2} & \textbf{Jaccard r2} & \textbf{Orden igual} & \textbf{Ruta completa} \\ \midrule
\endfirsthead
\toprule
\textbf{Protocolo} & \textbf{IP ida r2} & \textbf{IP vuelta r2} & \textbf{Jaccard r2} & \textbf{Orden igual} & \textbf{Ruta completa} \\ \midrule
\endhead
UDP & 9 & 7 & 0.071 & no & no \\
ICMP & 9 & 7 & 0.077 & no & sí \\
TCP & 9 & 6 & 0.077 & no & no \\
\bottomrule
\end{longtable}
\endgroup
\begingroup
\fontsize{9}{11}\selectfont
\begin{longtable}{p{0.95cm}p{7.075cm}p{7.075cm}}
\caption{Tabla tipo 2: Valparaíso / UDP, ronda 2.}\label{tabla2:138.204.230.242:UDP} \\
\toprule
\textbf{HOP} & \textbf{IP\_\allowbreak{}IDA} & \textbf{IP\_\allowbreak{}VUELTA} \\ \midrule
\endfirsthead
\toprule
\textbf{HOP} & \textbf{IP\_\allowbreak{}IDA} & \textbf{IP\_\allowbreak{}VUELTA} \\ \midrule
\endhead
1 & 192.168.100.1 [priv./esp.] & 138.204.230.193 [AS61506] \\
2 & 100.88.0.1 [priv./esp.] & 138.204.228.164 [AS61506] \\
3 & 10.87.0.3 [priv./esp.] & 10.100.163.1 [priv./esp.] \\
4 & 10.87.0.46 [priv./esp.] & 10.100.154.2 [priv./esp.] \\
5 & 45.68.16.129 [AS n/d] & 45.68.16.231 [AS n/d] \\
6 & 10.100.154.1 [priv./esp.] & 192.168.113.5 [priv./esp.] \\
7 & 138.204.228.164 [AS61506] & * \\
8 & 138.204.228.162 [AS61506] & * \\
9 & 138.204.230.242 [AS61506] & * \\
10 & -- & * \\
11 & -- & *; 138.204.228.163 [AS61506]; 138.204.230.193 [AS61506] \\
255 & -- & * \\
\bottomrule
\end{longtable}
\endgroup
\begingroup
\fontsize{9}{11}\selectfont
\begin{longtable}{p{0.95cm}p{7.075cm}p{7.075cm}}
\caption{Tabla tipo 2: Valparaíso / ICMP, ronda 2.}\label{tabla2:138.204.230.242:ICMP} \\
\toprule
\textbf{HOP} & \textbf{IP\_\allowbreak{}IDA} & \textbf{IP\_\allowbreak{}VUELTA} \\ \midrule
\endfirsthead
\toprule
\textbf{HOP} & \textbf{IP\_\allowbreak{}IDA} & \textbf{IP\_\allowbreak{}VUELTA} \\ \midrule
\endhead
1 & 192.168.100.1 [priv./esp.] & 138.204.230.193 [AS61506] \\
2 & 100.88.0.1 [priv./esp.] & 138.204.228.164 [AS61506] \\
3 & 10.87.0.2 [priv./esp.] & 10.100.163.1 [priv./esp.] \\
4 & 10.87.0.46 [priv./esp.] & 10.100.154.2 [priv./esp.] \\
5 & 45.68.16.129 [AS n/d] & 45.68.16.231 [AS n/d] \\
6 & 10.100.154.1 [priv./esp.] & 192.168.113.5 [priv./esp.] \\
7 & 138.204.228.164 [AS61506] & 186.189.89.115 [AS52341] \\
8 & 138.204.228.162 [AS61506] & -- \\
9 & 138.204.230.242 [AS61506] & -- \\
\bottomrule
\end{longtable}
\endgroup
\begingroup
\fontsize{9}{11}\selectfont
\begin{longtable}{p{0.95cm}p{7.075cm}p{7.075cm}}
\caption{Tabla tipo 2: Valparaíso / TCP, ronda 2.}\label{tabla2:138.204.230.242:TCP} \\
\toprule
\textbf{HOP} & \textbf{IP\_\allowbreak{}IDA} & \textbf{IP\_\allowbreak{}VUELTA} \\ \midrule
\endfirsthead
\toprule
\textbf{HOP} & \textbf{IP\_\allowbreak{}IDA} & \textbf{IP\_\allowbreak{}VUELTA} \\ \midrule
\endhead
1 & 192.168.100.1 [priv./esp.] & 138.204.230.193 [AS61506] \\
2 & 100.88.0.1 [priv./esp.] & 138.204.228.164 [AS61506] \\
3 & 10.87.0.3 [priv./esp.] & 10.100.163.1 [priv./esp.] \\
4 & 10.87.0.46 [priv./esp.] & 10.100.154.2 [priv./esp.] \\
5 & 45.68.16.129 [AS n/d] & 45.68.16.231 [AS n/d] \\
6 & 10.100.154.1 [priv./esp.] & 192.168.113.5 [priv./esp.] \\
7 & 138.204.228.164 [AS61506] & * \\
8 & 138.204.228.162 [AS61506] & * \\
9 & 138.204.230.242 [AS61506] & * \\
10 & -- & * \\
11 & -- & * \\
255 & -- & * \\
\bottomrule
\end{longtable}
\endgroup

En Valparaíso la vuelta es más corta que la ida, con ICMP se registran siete saltos hasta llegar a mi propia dirección (186.189.89.115, AS52341), contra nueve de la ida. Casi no hay direcciones en común (Jaccard entre 0.07 y 0.08), pero los segmentos se parecen, 45.68.16.129 en la ida y 45.68.16.231 en la vuelta, 10.100.154.1 y 10.100.154.2, y 138.204.228.164, que sí aparece en ambos sentidos (es lo que se esperaría de un camino que recorre los mismos enlaces en sentidos opuestos). Aun así la vuelta incluye una interfaz que la ida no muestra (10.100.163.1), por lo que el orden no resulta idéntico. Con UDP y TCP la vuelta se corta en el TTL 6 y no alcanza mi dirección, solo ICMP cierra la ruta. Es una ruta que se ve bastante simétrica, pero con huecos suficientes para no afirmarlo. Las tres rondas y las respuestas múltiples están en el anexo.


\subsubsection{La Habana}
\begingroup
\fontsize{9}{11}\selectfont
\addtocounter{table}{-1}
\begin{longtable}{p{2.1cm}p{1.7cm}p{2cm}p{2cm}p{2.2cm}p{2.7cm}}
\toprule
\textbf{Protocolo} & \textbf{IP ida r2} & \textbf{IP vuelta r2} & \textbf{Jaccard r2} & \textbf{Orden igual} & \textbf{Ruta completa} \\ \midrule
\endfirsthead
\toprule
\textbf{Protocolo} & \textbf{IP ida r2} & \textbf{IP vuelta r2} & \textbf{Jaccard r2} & \textbf{Orden igual} & \textbf{Ruta completa} \\ \midrule
\endhead
UDP & 15 & 11 & 0.000 & no & no \\
ICMP & 13 & 12 & 0.000 & no & no \\
TCP & 16 & 10 & 0.000 & no & no \\
\bottomrule
\end{longtable}
\endgroup
\begingroup
\fontsize{9}{11}\selectfont
\begin{longtable}{p{0.95cm}p{7.075cm}p{7.075cm}}
\caption{Tabla tipo 2: La Habana / UDP, ronda 2.}\label{tabla2:200.14.48.134:UDP} \\
\toprule
\textbf{HOP} & \textbf{IP\_\allowbreak{}IDA} & \textbf{IP\_\allowbreak{}VUELTA} \\ \midrule
\endfirsthead
\toprule
\textbf{HOP} & \textbf{IP\_\allowbreak{}IDA} & \textbf{IP\_\allowbreak{}VUELTA} \\ \midrule
\endhead
1 & 192.168.100.1 [priv./esp.] & 200.14.48.132 [AS27725] \\
2 & 100.88.0.1 [priv./esp.] & 200.14.48.177 [AS27725] \\
3 & 10.87.0.3 [priv./esp.] & 192.168.118.89 [priv./esp.] \\
4 & 10.87.0.46 [priv./esp.] & 172.31.247.241 [priv./esp.] \\
5 & * & 200.0.16.30 [AS11960] \\
6 & 154.54.92.41 [AS174] & 200.0.16.18 [AS11960] \\
7 & *; 154.54.29.157 [AS174] & 200.0.16.20 [AS11960] \\
8 & 154.54.167.185 [AS174] & 89.221.37.138 [AS6762] \\
9 & 38.88.57.26 [AS174] & 8.243.236.5 [AS3356] \\
10 & * & 200.189.207.14 [AS3549] \\
11 & 69.79.104.1 [AS23520] & 8.243.191.90 [AS3356] \\
12 & * & * \\
13 & 190.242.168.39 [AS23520] & * \\
14 & 200.0.16.16 [AS11960] & *; 200.14.48.132 [AS27725]; 200.14.48.177 [AS27725] \\
15 & 200.0.16.21 [AS11960] & * \\
16 & 200.0.16.29 [AS11960] & * \\
17 & * & -- \\
18 & * & -- \\
19 & 200.14.48.178 [AS27725] & -- \\
20 & 200.14.48.134 [AS27725] & -- \\
255 & -- & * \\
\bottomrule
\end{longtable}
\endgroup
\begingroup
\fontsize{9}{11}\selectfont
\begin{longtable}{p{0.95cm}p{7.075cm}p{7.075cm}}
\caption{Tabla tipo 2: La Habana / ICMP, ronda 2.}\label{tabla2:200.14.48.134:ICMP} \\
\toprule
\textbf{HOP} & \textbf{IP\_\allowbreak{}IDA} & \textbf{IP\_\allowbreak{}VUELTA} \\ \midrule
\endfirsthead
\toprule
\textbf{HOP} & \textbf{IP\_\allowbreak{}IDA} & \textbf{IP\_\allowbreak{}VUELTA} \\ \midrule
\endhead
1 & 192.168.100.1 [priv./esp.] & 200.14.48.132 [AS27725] \\
2 & 100.88.0.1 [priv./esp.] & 200.14.48.177 [AS27725] \\
3 & 10.87.0.2 [priv./esp.] & 192.168.118.89 [priv./esp.] \\
4 & 10.87.0.46 [priv./esp.] & 172.31.247.241 [priv./esp.] \\
5 & * & 200.0.16.86 [AS11960] \\
6 & 154.54.92.41 [AS174] & 200.0.16.45 [AS11960] \\
7 & 154.54.29.157 [AS174] & 200.0.16.20 [AS11960] \\
8 & 154.54.167.185 [AS174] & 89.221.37.138 [AS6762] \\
9 & 38.104.94.178 [AS174] & 8.243.236.5 [AS3356] \\
10 & * & 200.189.207.14 [AS3549] \\
11 & 190.242.168.37 [AS23520] & 8.243.191.90 [AS3356] \\
12 & 200.0.16.21 [AS11960] & *; 186.189.89.115 [AS52341] \\
13 & 200.0.16.29 [AS11960] & -- \\
14 & * & -- \\
15 & * & -- \\
16 & 200.14.48.178 [AS27725] & -- \\
17 & 200.14.48.134 [AS27725] & -- \\
\bottomrule
\end{longtable}
\endgroup
\begingroup
\fontsize{9}{11}\selectfont
\begin{longtable}{p{0.95cm}p{7.075cm}p{7.075cm}}
\caption{Tabla tipo 2: La Habana / TCP, ronda 2.}\label{tabla2:200.14.48.134:TCP} \\
\toprule
\textbf{HOP} & \textbf{IP\_\allowbreak{}IDA} & \textbf{IP\_\allowbreak{}VUELTA} \\ \midrule
\endfirsthead
\toprule
\textbf{HOP} & \textbf{IP\_\allowbreak{}IDA} & \textbf{IP\_\allowbreak{}VUELTA} \\ \midrule
\endhead
1 & 192.168.100.1 [priv./esp.] & 200.14.48.132 [AS27725] \\
2 & 100.88.0.1 [priv./esp.] & 200.14.48.177 [AS27725] \\
3 & 10.87.0.3 [priv./esp.] & 192.168.118.89 [priv./esp.] \\
4 & 10.87.0.46 [priv./esp.] & *; 200.14.48.132 [AS27725]; 200.14.48.177 [AS27725] \\
5 & * & 200.0.16.86 [AS11960] \\
6 & 154.54.92.41 [AS174] & 200.0.16.45 [AS11960] \\
7 & *; 154.54.29.157 [AS174]; 154.54.3.29 [AS174] & 200.0.16.20 [AS11960] \\
8 & 154.54.167.185 [AS174] & 89.221.37.138 [AS6762] \\
9 & 38.88.57.26 [AS174] & 8.243.236.9 [AS3356] \\
10 & * & 200.189.207.14 [AS3549] \\
11 & *; 190.242.168.39 [AS23520] & 8.243.191.90 [AS3356] \\
12 & 200.0.16.16 [AS11960] & * \\
13 & 200.0.16.21 [AS11960] & * \\
14 & 200.0.16.25 [AS11960] & * \\
15 & *; 200.0.16.25 [AS11960] & * \\
16 & *; 192.168.100.1 [priv./esp.] & * \\
17 & *; 100.88.0.1 [priv./esp.]; 200.0.16.29 [AS11960] & -- \\
18 & * & -- \\
19 & 200.14.48.178 [AS27725] & -- \\
20 & 200.14.48.134 [AS27725] & -- \\
255 & -- & * \\
\bottomrule
\end{longtable}
\endgroup

Con la Habana la situación cambia, no hay ninguna IP en común entre ida y vuelta en ningún protocolo (Jaccard de 0.000) y tampoco coinciden los sistemas autónomos. La ida sale por AS174 y pasa por AS23520, AS11960 y AS27725. la vuelta sale de AS27725, atraviesa AS11960 y luego AS6762, AS3356 y AS3549 (200.189.207.14), y llega a mi conexión por una interfaz de AS3356 (8.243.191.90). Con ICMP la vuelta alcanza mi dirección, aunque con una respuesta perdida en el último TTL.
Que cambien los sistemas autónomos y no solo las direcciones, es una señal más fuerte de asimetría que un Jaccard bajo, porque el efecto de ver cada interfaz de entrada no alcanza a explicar un cambio de AS. Con todo, los huecos de ambos sentidos impiden decir que las rutas completas sean distintas, lo que se observa es que los tramos visibles no coinciden. Las tres rondas y las respuestas múltiples están en el anexo.


\subsubsection{Rusia}
\begingroup
\fontsize{9}{11}\selectfont
\addtocounter{table}{-1}
\begin{longtable}{p{2.1cm}p{1.7cm}p{2cm}p{2cm}p{2.2cm}p{2.7cm}}
\toprule
\textbf{Protocolo} & \textbf{IP ida r2} & \textbf{IP vuelta r2} & \textbf{Jaccard r2} & \textbf{Orden igual} & \textbf{Ruta completa} \\ \midrule
\endfirsthead
\toprule
\textbf{Protocolo} & \textbf{IP ida r2} & \textbf{IP vuelta r2} & \textbf{Jaccard r2} & \textbf{Orden igual} & \textbf{Ruta completa} \\ \midrule
\endhead
UDP & 11 & 8 & 0.000 & no & no \\
ICMP & 11 & 8 & 0.000 & no & no \\
TCP & 11 & 7 & 0.000 & no & no \\
\bottomrule
\end{longtable}
\endgroup
\begingroup
\fontsize{9}{11}\selectfont
\begin{longtable}{p{0.95cm}p{7.075cm}p{7.075cm}}
\caption{Tabla tipo 2: Rusia / UDP, ronda 2.}\label{tabla2:37.79.226.46:UDP} \\
\toprule
\textbf{HOP} & \textbf{IP\_\allowbreak{}IDA} & \textbf{IP\_\allowbreak{}VUELTA} \\ \midrule
\endfirsthead
\toprule
\textbf{HOP} & \textbf{IP\_\allowbreak{}IDA} & \textbf{IP\_\allowbreak{}VUELTA} \\ \midrule
\endhead
1 & 192.168.100.1 [priv./esp.] & 192.168.30.1 [priv./esp.] \\
2 & 100.88.0.1 [priv./esp.] & 37.79.226.33 [AS12389] \\
3 & 10.87.0.3 [priv./esp.] & 2.63.212.245 [AS12389] \\
4 & 10.87.0.46 [priv./esp.] & 172.29.8.149 [priv./esp.] \\
5 & * & *; 172.29.39.155 [priv./esp.] \\
6 & 154.54.92.41 [AS174] & 212.133.1.53 [AS3356] \\
7 & *; 154.54.29.157 [AS174] & 200.189.207.14 [AS3549] \\
8 & 154.54.167.185 [AS174] & 8.243.191.90 [AS3356] \\
9 & 154.54.10.210 [AS174] & * \\
10 & 195.22.211.61 [AS6762] & * \\
11 & *; 100.88.0.1 [priv./esp.]; 192.168.100.1 [priv./esp.]; 195.22.214.161 [AS6762] & * \\
12 & * & *; 192.168.30.1 [priv./esp.]; 37.79.226.33 [AS12389] \\
13 & * & * \\
14 & 2.63.212.246 [AS12389] & -- \\
15 & * & -- \\
16 & * & -- \\
17 & * & -- \\
18 & * & -- \\
19 & * & -- \\
255 & * & * \\
\bottomrule
\end{longtable}
\endgroup
\begingroup
\fontsize{9}{11}\selectfont
\begin{longtable}{p{0.95cm}p{7.075cm}p{7.075cm}}
\caption{Tabla tipo 2: Rusia / ICMP, ronda 2.}\label{tabla2:37.79.226.46:ICMP} \\
\toprule
\textbf{HOP} & \textbf{IP\_\allowbreak{}IDA} & \textbf{IP\_\allowbreak{}VUELTA} \\ \midrule
\endfirsthead
\toprule
\textbf{HOP} & \textbf{IP\_\allowbreak{}IDA} & \textbf{IP\_\allowbreak{}VUELTA} \\ \midrule
\endhead
1 & 192.168.100.1 [priv./esp.] & 192.168.30.1 [priv./esp.] \\
2 & 100.88.0.1 [priv./esp.] & 37.79.226.33 [AS12389] \\
3 & 10.87.0.2 [priv./esp.] & 2.63.212.245 [AS12389] \\
4 & 10.87.0.46 [priv./esp.] & 172.29.8.149 [priv./esp.] \\
5 & * & * \\
6 & 154.54.92.41 [AS174] & 212.133.1.53 [AS3356] \\
7 & *; 154.54.3.29 [AS174] & 200.189.207.14 [AS3549] \\
8 & 154.54.168.85 [AS174] & 8.243.191.90 [AS3356] \\
9 & 154.54.10.210 [AS174] & 186.189.89.115 [AS52341] \\
10 & 195.22.211.61 [AS6762] & -- \\
11 & * & -- \\
12 & * & -- \\
13 & * & -- \\
14 & 2.63.212.246 [AS12389] & -- \\
15 & 37.79.226.46 [AS12389] & -- \\
\bottomrule
\end{longtable}
\endgroup
\begingroup
\fontsize{9}{11}\selectfont
\begin{longtable}{p{0.95cm}p{7.075cm}p{7.075cm}}
\caption{Tabla tipo 2: Rusia / TCP, ronda 2.}\label{tabla2:37.79.226.46:TCP} \\
\toprule
\textbf{HOP} & \textbf{IP\_\allowbreak{}IDA} & \textbf{IP\_\allowbreak{}VUELTA} \\ \midrule
\endfirsthead
\toprule
\textbf{HOP} & \textbf{IP\_\allowbreak{}IDA} & \textbf{IP\_\allowbreak{}VUELTA} \\ \midrule
\endhead
1 & 192.168.100.1 [priv./esp.] & 192.168.30.1 [priv./esp.] \\
2 & 100.88.0.1 [priv./esp.] & 37.79.226.33 [AS12389] \\
3 & 10.87.0.3 [priv./esp.] & 2.63.212.245 [AS12389] \\
4 & 10.87.0.46 [priv./esp.] & 172.29.8.149 [priv./esp.] \\
5 & * & * \\
6 & *; 154.54.92.41 [AS174] & 212.133.1.53 [AS3356] \\
7 & *; 154.54.3.29 [AS174] & 200.189.207.14 [AS3549] \\
8 & 154.54.168.85 [AS174] & 8.243.191.90 [AS3356] \\
9 & 154.54.10.210 [AS174] & * \\
10 & 195.22.211.209 [AS6762] & * \\
11 & * & * \\
12 & * & * \\
13 & * & * \\
14 & *; 2.63.212.246 [AS12389] & -- \\
15 & 37.79.226.46 [AS12389] & -- \\
16 & 37.79.226.46 [AS12389] & -- \\
17 & *; 37.79.226.46 [AS12389] & -- \\
255 & -- & * \\
\bottomrule
\end{longtable}
\endgroup

En Rusia la imagen es parecida, Jaccard de 0.000, ninguna IP en común y ningún ASN coincidente en lo visible. La ida entra por AS174 y AS6762 hasta AS12389, la vuelta sale directamente desde AS12389 (37.79.226.33 y 2.63.212.245) hacia AS3356 (212.133.1.53),  AS3549 (200.189.207.14) y AS3356 (8.243.191.90). Estas dos últimas interfaces son las mismas que aparecen en la vuelta desde Cuba justo antes de mi dirección, el tráfico que vuelve a mi casa desde dos países distintos converge en las mismas interfaces de entrada a la red del operador, mientras que la ida a esos dos destinos compartía el comienzo por AS174. Con ICMP la vuelta alcanza mi dirección en el noveno salto, con UDP y TCP se pierde antes. Las tres rondas y las respuestas están en el anexo.


